\chapter{Choix techniques}

\section{Architecture en couches}

Chaque requête traverse les couches dans un seul sens : filtre → servlet → service → DAO → Hibernate.

\begin{itemize}
  \item \textbf{Servlets} (\code{controller}) : lecture des paramètres, appel d'un service, dépôt des attributs, puis
        \code{forward} vers une JSP de \code{WEB-INF/views} ou redirection après un POST. Aucune règle métier ni requête HQL.
  \item \textbf{Services} (\code{service}) : toutes les règles de gestion et les transactions. Une violation de règle lève
        une \code{RegleMetierException} dont le message, en français, est affiché par la JSP.
  \item \textbf{DAO} (\code{dao}) : accès aux données uniquement. \code{GenericDao<T>} fournit le CRUD commun (\code{sauvegarder},
        \code{modifier}, \code{supprimer}, \code{trouverParId}, \code{listerTous}), les DAO spécifiques ajoutent les requêtes HQL.
  \item \textbf{Modèle} (\code{model}) : entités JPA annotées et objets du panier.
  \item \textbf{Utilitaires} (\code{util}) : \code{HibernateUtil}, \code{DemarrageListener} (création de la \code{SessionFactory} au
        déploiement, fermeture à l'arrêt), \code{Parametres} (lecture des paramètres numériques : une valeur absente ou
        mal formée ne provoque jamais d'erreur 500).
\end{itemize}

\section{Hibernate natif et session liée au thread}

\begin{itemize}
  \item \code{SessionFactory} unique créée par \code{HibernateUtil} à partir de \code{hibernate.cfg.xml}, sans couche
        \code{EntityManager} ni framework d'injection.
  \item \code{hibernate.current\_session\_context\_class = thread} : le service ouvre la transaction sur
        \code{getCurrentSession()}, les DAO utilisent la même session sans la gérer. Chaque service dispose d'une
        méthode privée \code{enTransaction} qui valide (commit) ou annule (rollback) la transaction.
  \item Tables générées par \code{hbm2ddl.auto = update} à partir des annotations.
  \item Associations \code{@ManyToOne} chargées immédiatement (comportement JPA par défaut), ce qui permet aux JSP
        d'afficher la catégorie d'un produit ou le client d'une commande après la fermeture de la session.
        \code{Commande.lignes} est chargée à la demande ; le détail d'une commande utilise une requête \code{join fetch}
        (\code{CommandeDao.trouverAvecLignes}) pour charger lignes et produits en une seule fois.
\end{itemize}

\section{Cohérence du stock}

\begin{itemize}
  \item \textbf{Passage de commande} : chaque produit est relu avec un verrou pessimiste en écriture
        (\code{LockMode.PESSIMISTIC\_WRITE}, soit \code{select … for update}). Deux commandes simultanées sur le même produit
        sont traitées l'une après l'autre, ce qui empêche de vendre plus que le stock. La requête HQL ne verrouille
        que la ligne du produit (et non celle de sa catégorie), et les produits sont verrouillés par identifiant
        croissant afin que deux transactions prennent les verrous dans le même ordre et ne puissent pas se bloquer
        mutuellement.
  \item \textbf{Annulation} : la commande est verrouillée avant le changement de statut, si bien que deux annulations
        simultanées ne réintègrent le stock qu'une fois. Les quantités sont rendues par une requête HQL
        \code{update Produit p set p.stock = p.stock + :quantite}, calculée par la base sur la valeur actuelle du stock :
        une vente validée entre-temps par une autre transaction n'est pas écrasée.
  \item \textbf{Prix} : le prix est relu en base à la validation et copié dans la ligne de commande (RG10) ; le prix
        conservé dans le panier ne sert qu'à l'affichage.
\end{itemize}

\section{Panier en session}

Le panier (\code{Panier}, \code{LignePanier}) est un objet simple stocké dans la \code{HttpSession}, jamais persisté
(cahier des charges §7). Les lignes sont indexées par identifiant de produit dans une \code{LinkedHashMap} (ordre d'ajout
conservé). Le service recharge le produit à chaque ajout ou modification pour contrôler le stock réel.
Conséquence assumée : le panier est perdu à la déconnexion ou à l'expiration de la session.

\section{Navigation : PRG et messages flash}

Après chaque POST qui modifie des données, la servlet redirige (modèle Post/Redirect/Get) : un
rafraîchissement de la page ne renvoie pas le formulaire. Le message de confirmation ou d'erreur est déposé
en session (\code{flashSucces}, \code{flashErreur}) puis affiché une seule fois et retiré par l'en-tête commun.
Quand la saisie doit être corrigée (inscription, profil, validation de commande), la servlet fait un
\code{forward} vers le formulaire avec le message et les valeurs saisies.

\section{Sécurité}\label{sec:securite}

\begin{itemize}
  \item \textbf{Contrôle d'accès} par filtres (BNF02) : \code{AuthFilter} (\code{/client/*}) exige le rôle CLIENT, un visiteur
        est redirigé vers la connexion et un administrateur vers son tableau de bord ; \code{AdminFilter} (\code{/admin/*})
        exige le rôle ADMIN, un visiteur est redirigé vers la connexion et un client reçoit une erreur 403.
        Les trois filtres (\code{EncodingFilter}, \code{AuthFilter}, \code{AdminFilter}) sont déclarés dans \code{web.xml} : seul
        l'ordre des \code{<filter-mapping>} est garanti par la spécification Servlet, ce qui assure que
        \code{EncodingFilter} s'exécute en premier.
  \item \textbf{Mots de passe} hachés avec BCrypt (sel intégré) ; message d'échec identique que l'e-mail existe ou non.
  \item \textbf{Injection SQL} : toutes les requêtes HQL utilisent des paramètres nommés ; la requête de recherche du
        catalogue est construite dynamiquement, mais seules des clauses fixes sont concaténées, jamais une valeur
        saisie.
  \item \textbf{XSS} : toutes les données saisies par les utilisateurs sont affichées par \code{<c:out>}, qui échappe le HTML.
  \item \textbf{Session} : nouvel identifiant de session à la connexion (contre la fixation de session), identifiant
        transmis uniquement par cookie (\code{tracking-mode COOKIE}, jamais dans l'URL), invalidation à la déconnexion,
        expiration après 30 minutes. La déconnexion se fait uniquement en POST.
  \item \textbf{Contrôles côté serveur} : toutes les règles sont vérifiées dans les services, même si un formulaire
        est contourné ; par exemple, un statut de commande forgé est refusé et une commande d'un autre client est
        traitée comme inexistante (RG16).
\end{itemize}

\section{Présentation}

\begin{itemize}
  \item JSP uniquement en JSTL et EL, sans scriptlet (BNF06) ; en-tête et pied de page communs inclus dans chaque
        page, fragments communs pour la pagination, le badge de statut et le résumé d'une commande.
  \item Bootstrap 5.3 par CDN, mise en page responsive (BNF08).
  \item Prix affichés avec \code{fmt:formatNumber} (deux décimales, format français, en DH) ; dates formatées par
        l'entité (\code{getDateCommandeFormatee}), car \code{fmt:formatDate} ne prend pas en charge \code{LocalDateTime}.
  \item Libellés français des statuts portés par l'énumération \code{StatutCommande}.
\end{itemize}

\section{Pages d'erreur}

Une seule page \code{commun/erreur.jsp}, déclarée dans \code{web.xml} pour les codes 400, 403, 404, 405 et 500 et pour
toute exception non interceptée. Elle affiche un message en français avec l'en-tête et le pied communs, sans
trace technique ni version du serveur.

\section{UTF-8 de bout en bout}

Base et tables en \code{utf8mb4} (collation \code{utf8mb4\_unicode\_ci}, qui rend aussi la recherche insensible aux
accents), connexion JDBC en UTF-8, \code{EncodingFilter} sur toutes les requêtes (lecture des paramètres et
réponse), encodage des JSP fixé dans \code{web.xml}, sources compilées en UTF-8 (\code{pom.xml}).
